% ============================================================================
% Chapter 14: Coordination Frustration Correction
% ============================================================================
% Source: radioactivity_v7_8_Salpha_deformation, superheavy_predictions.py
% Status: FILLED from SoT
% Contamination check: PASSED (no banned terms)
% Contamination guard: SM/nuclear-model terminology routed to Appendix X/Q.
% Layer A text uses EDC-native vocabulary only.
% ============================================================================

\chapter{Coordination Frustration Correction}
\label{ch:frustrated}

\source{radioactivity\_v7\_8\_Salpha\_deformation, superheavy\_predictions.py}

% ----------------------------------------------------------------------------
\section*{Abstract}
% ----------------------------------------------------------------------------

This chapter develops the \textbf{coordination frustration correction} that
explains baseline residuals from Chapter~\ref{ch:barrier_release}. Starting
from the residual definition $\Delta := \log_{10}(t_{\text{obs}}/t_{\text{BL}})$,
we introduce the coordination function $n(A)$, define the allowed set
$S = \{2^a \times 3^b\}$ from M$_6$ topology, and construct the distance-to-allowed
metric $d(n)$. The corrected lane becomes
$\Delta \approx g \times d(n)$, where the coupling coefficient $g < 0$ indicates
that higher frustration leads to \emph{faster} release. This chapter provides
the theoretical framework; calibration of $g$ is routed to Appendix~Q.

% ----------------------------------------------------------------------------
% CHAPTER SPINE
% ----------------------------------------------------------------------------
\paragraph{Chapter Spine.}
\textbf{Purpose:} Derive the frustration correction $\Delta = g \times d(n)$ that
explains baseline residuals using $\Msix$ topology.
\textbf{Inputs:} (i)~Baseline residuals from Ch.~\ref{ch:barrier_release}~[BL],
(ii)~allowed set $S = \{2^a \times 3^b\}$ from Ch.~\ref{ch:M6}~[Der],
(iii)~coordination function $n(A)$~[Cal].
\textbf{Outputs:} (i)~Distance-to-allowed $d(n) = \min_{k \in S}|n-k|$~[Def],
(ii)~frustration correction $\Delta = g \times d(n)$~[Dc],
(iii)~$g$ coupling coefficient~[I].
\textbf{Dependencies:} Ch.~\ref{ch:M6} (allowed set $S$), Ch.~\ref{ch:barrier_release} (residuals).
\textbf{Forward links:} Correction applied in Ch.~\ref{ch:high_coord} (superheavy predictions).

% ============================================================================
\section{From Baseline Residuals to Structural Correction}
\label{sec:frustrated:motivation}
% ============================================================================

\subsection{The Residual Interface from Chapter~\ref{ch:barrier_release}}

Chapter~\ref{ch:barrier_release} established the baseline lane for barrier-limited
closed-4 emission:
\begin{equation}
    \log_{10}(t_{1/2}) = a \times X + b \tagBL
    \label{eq:f:baseline_recall}
\end{equation}
where $X = Z_{\text{eff}}/\sqrt{E_{\text{rel}}}$ is the barrier-to-release ratio.

The baseline explains $\gtrsim 98\%$ of variance in $\log_{10}(t_{1/2})$ across
$\sim 20$ orders of magnitude. However, systematic residuals remain:

\begin{derivationbox}
    \textbf{Baseline residual (from \S\ref{sec:gn:residual}):}
    \begin{equation}
        \Delta := \log_{10}(t_{\text{obs}}) - \log_{10}(t_{\text{BL}})
        = \log_{10}\!\left(\frac{t_{\text{obs}}}{t_{\text{BL}}}\right)
        \label{eq:f:residual_def}
    \end{equation}
    \tagBL
\end{derivationbox}

\subsection{Why Residuals Encode Structure}

The baseline treats the parent cluster as featureless. But parent clusters
are \textbf{junction networks} (Chapter~\ref{ch:deuterium}) with specific
coordination structures. When the coordination number $n$ of a cluster
deviates from the allowed set $S$ determined by M$_6$ topology
(Chapter~\ref{ch:M6}), structural strain modifies the effective barrier.

\begin{resultbox}
    \textbf{Central hypothesis:}
    Baseline residuals $\Delta$ correlate with a measurable \emph{coordination
    frustration metric} $d(n)$ derived from the cluster's position relative to
    the allowed coordination set.
    \tagP
\end{resultbox}

% ============================================================================
\section{The Coordination Function $n(A)$}
\label{sec:frustrated:nA}
% ============================================================================

\subsection{Definition}

The \textbf{effective coordination number} of a parent cluster with mass number
$A$ is defined by the scaling relation:

\begin{derivationbox}
    \begin{equation}
        n(A) = p \times A^{1/3}
        \label{eq:f:nA}
    \end{equation}
    where $p$ is a prefactor calibrated to known junction networks.
    \tagCal
\end{derivationbox}

The $A^{1/3}$ scaling arises from surface-to-volume considerations: the
number of surface coordination sites scales with the cluster's linear
dimension $\sim A^{1/3}$.

\subsection{Prefactor Calibration}

\begin{quarantinebox}[title=Appendix Q: Prefactor Calibration]
    The prefactor $p$ is calibrated so that $n(208) \approx 36$ for the
    doubly-stable Pb-208 cluster. This places the heaviest stable
    configuration at the boundary of the allowed set.
    Numerical value: $p \approx 6.1$. Full calibration procedure in Appendix~Q.
\end{quarantinebox}

\subsection{EDC-Native Interpretation}

In EDC terms, $n(A)$ measures:
\begin{itemize}
    \item The \textbf{effective coordination demand} of the parent junction network
    \item The number of boundary junctions relevant for closed-4 tunneling
    \item The topological complexity of the parent configuration
\end{itemize}

This is \emph{not} a conventional count of constituents, but a
topological measure of coordination structure. \tagP

% ============================================================================
\section{The Allowed Set $S = \{2^a \times 3^b\}$}
\label{sec:frustrated:allowed}
% ============================================================================

\subsection{Definition and Origin}

The \textbf{allowed coordination set} consists of integers expressible as
products of powers of 2 and 3:

\begin{derivationbox}
    \begin{equation}
        S = \{n : n = 2^a \times 3^b, \quad a, b \geq 0\}
        \label{eq:f:allowed_set}
    \end{equation}
    \tagDer
\end{derivationbox}

This set arises from the Z$_6$ symmetry of the junction topology
(Chapter~\ref{ch:junction}). Since $\text{Z}_6 \cong \text{Z}_2 \times \text{Z}_3$,
allowed coordination numbers must be compatible with both the binary and
ternary symmetry factors.

\subsection{Explicit Enumeration}

Table~\ref{tab:f:allowed} lists the first 23 allowed coordination numbers.

\begin{table}[h]
\centering
\caption{Allowed coordination numbers $S = \{2^a \times 3^b\}$.}
\label{tab:f:allowed}
\begin{tabular}{cccc}
    \toprule
    $n$ & $(a, b)$ & $n$ & $(a, b)$ \\
    \midrule
    1 & (0,0) & 24 & (3,1) \\
    2 & (1,0) & 27 & (0,3) \\
    3 & (0,1) & 32 & (5,0) \\
    4 & (2,0) & 36 & (2,2) \\
    6 & (1,1) & 48 & (4,1) \\
    8 & (3,0) & 54 & (1,3) \\
    9 & (0,2) & 64 & (6,0) \\
    12 & (2,1) & 72 & (3,2) \\
    16 & (4,0) & 81 & (0,4) \\
    18 & (1,2) & 96 & (5,1) \\
    & & 108 & (2,3) \\
    \bottomrule
\end{tabular}
\end{table}

\subsection{The Forbidden Zone}

A striking feature of the allowed set is the existence of \textbf{forbidden
zones}---intervals containing no allowed values:

\begin{resultbox}
    \textbf{Primary forbidden zone:}
    \begin{equation}
        [37, 47] \subset \mathbb{Z}^+ \setminus S
    \end{equation}
    All 11 integers in this range are forbidden by M$_6$ topology.
    \tagDer
\end{resultbox}

Clusters with $n(A)$ falling in this zone experience maximal coordination
frustration. This has dramatic consequences for high-coordination predictions
(Chapter~\ref{ch:high_coord}).

\subsection{Gap Structure}

The gaps between consecutive allowed values grow with $n$:

\begin{center}
\begin{tabular}{ccc}
    \toprule
    Allowed pair & Gap & Notes \\
    \midrule
    $36 \to 48$ & 12 & Contains forbidden zone [37--47] \\
    $48 \to 54$ & 6 & \\
    $54 \to 64$ & 10 & \\
    $64 \to 72$ & 8 & \\
    $72 \to 81$ & 9 & \\
    \bottomrule
\end{tabular}
\end{center}

The widening gaps imply that heavy clusters are increasingly likely to
experience significant coordination frustration. \tagP

% ============================================================================
\section{Distance-to-Allowed: The $d(n)$ Metric}
\label{sec:frustrated:dn}
% ============================================================================

\subsection{Definition}

The \textbf{coordination frustration distance} measures how far a cluster's
effective coordination $n(A)$ lies from the nearest allowed value:

\begin{derivationbox}
    \begin{equation}
        d(n) = \min_{k \in S} |n(A) - k|
        \label{eq:f:d_n}
    \end{equation}
    where $S = \{2^a \times 3^b\}$ is the allowed set.
    \tagDer
\end{derivationbox}

\subsection{Properties}

\begin{itemize}
    \item $d(n) = 0$ when $n(A)$ coincides with an allowed value
    \item $d(n) \in [0, 5.5]$ in practice (maximum when $n(A) \approx 42.5$
          in the forbidden zone center)
    \item $d(n)$ is piecewise linear between consecutive allowed values
    \item Larger $d(n)$ indicates greater topological mismatch
\end{itemize}

\subsection{Example Calculations}

Table~\ref{tab:f:examples} shows example $d(n)$ calculations for representative
clusters.

\begin{table}[h]
\centering
\caption{Example coordination frustration calculations.}
\label{tab:f:examples}
\begin{tabular}{lccccc}
    \toprule
    Cluster & $A$ & $n(A)$ & Nearest allowed & $d(n)$ & Zone \\
    \midrule
    Pb-208 & 208 & 36.1 & 36 & 0.1 & Allowed \\
    U-238 & 238 & 37.8 & 36 & 1.8 & Forbidden \\
    Cm-248 & 248 & 38.3 & 36 & 2.3 & Forbidden \\
    Fm-256 & 256 & 38.7 & 36 & 2.7 & Forbidden \\
    Og-294 & 294 & 40.5 & 36 & 4.5 & Deep forbidden \\
    \bottomrule
\end{tabular}
\end{table}

\noindent
Note: Heavier clusters fall deeper into the forbidden zone. \tagCal

% ============================================================================
\section{The Corrected Lane: $\Delta \approx g \times d(n)$}
\label{sec:frustrated:corrected}
% ============================================================================

\subsection{Form of the Correction}

The EDC correction to the baseline lane takes the linear form:

\begin{derivationbox}
    \textbf{Corrected release time:}
    \begin{equation}
        \log_{10}(t_{1/2}) = a X + b + g \times d(n)
        \label{eq:f:corrected_lane}
    \end{equation}
    or equivalently:
    \begin{equation}
        \Delta \approx g \times d(n)
        \label{eq:f:delta_g}
    \end{equation}
    where the coupling coefficient $g$ is identified \tagI{} from the
    $\Msix$ frustration structure; numerical calibration in Appendix~Q.
    \tagP
\end{derivationbox}

\subsection{The Sign of $g$}

The coupling coefficient $g$ \tagI{} is identified from the theoretical framework
and calibrated from data (Appendix~Q). Empirically:

\begin{resultbox}
    \textbf{Sign determination:}
    \begin{equation}
        g < 0
    \end{equation}
    Higher coordination frustration leads to \emph{faster} release
    (shorter half-life).
    \tagCal
\end{resultbox}

\subsection{Physical Interpretation}

The sign $g < 0$ admits two complementary interpretations:

\begin{enumerate}
    \item \textbf{Enhanced tunneling:} Frustrated configurations have higher
          internal energy, reducing the effective barrier height
    \item \textbf{Enhanced preformation:} Frustrated networks more readily
          form mobile closed-4 units that can escape
\end{enumerate}

Both interpretations are consistent with the topological frustration picture:
configurations that cannot achieve allowed coordination are unstable and
release closed-4 units more readily. \tagP

\subsection{Robustness}

\begin{quarantinebox}[title=Appendix Q: Robustness Analysis]
    The coefficient $g$ has been tested against multiple confounding variables:
    \begin{itemize}
        \item Boundary mismatch proxies: $g$ survives (4\% change)
        \item Preformation proxies: $g$ survives (7\% change)
        \item Prefactor variation $p \pm 0.1$: $g$ stable within uncertainties
    \end{itemize}
    Full robustness tables in Appendix~Q.
\end{quarantinebox}

% ============================================================================
\section{Model Performance}
\label{sec:frustrated:performance}
% ============================================================================

\subsection{Variance Reduction}

The coordination frustration correction improves baseline predictions:

\begin{center}
\begin{tabular}{lcc}
    \toprule
    Model & Mean $|\Delta|$ (dex) & Status \\
    \midrule
    Baseline only & $\sim 1.0$ & [BL] \\
    Baseline $+$ $g \times d(n)$ & $\sim 0.3$ & [BL] + [P] \\
    \bottomrule
\end{tabular}
\end{center}

The correction reduces mean absolute residual by approximately 3$\times$.

\subsection{High-Coordination Extrapolation}

For the high-A region (large clusters), where configurations fall deep in the
forbidden zone, the correction is essential:

\begin{center}
\begin{tabular}{lcc}
    \toprule
    Model & Mean $|\Delta|$ (high-A) & Factor \\
    \midrule
    Baseline only & $\sim 6$ dex & --- \\
    Baseline $+$ $g \times d(n)$ & $\sim 1$ dex & 6$\times$ \\
    \bottomrule
\end{tabular}
\end{center}

Without the frustration correction, baseline predictions for Og-294
are wrong by $\sim 6$ orders of magnitude (weeks instead of milliseconds).
\tagCal

% ============================================================================
\section{Tables}
\label{sec:frustrated:tables}
% ============================================================================

\subsection{Table 14.1: Notation Dictionary}

\begin{table}[h]
\centering
\caption{Notation dictionary for coordination frustration.}
\label{tab:f:notation}
\begin{tabular}{lp{7cm}c}
    \toprule
    Symbol & Definition & Tag \\
    \midrule
    $n(A)$ & Effective coordination number: $p \times A^{1/3}$ & [Cal] \\
    $p$ & Prefactor ($\approx 6.1$) calibrated to Pb-208 & [Cal] \\
    $S$ & Allowed set: $\{2^a \times 3^b\}$ & [Der] \\
    $d(n)$ & Distance-to-allowed: $\min_{k \in S}|n - k|$ & [Der] \\
    $g$ & Frustration coupling coefficient & [Cal] \\
    $\Delta$ & Baseline residual: $\log_{10}(t_{\text{obs}}/t_{\text{BL}})$ & [BL] \\
    \bottomrule
\end{tabular}
\end{table}

\subsection{Table 14.2: Allowed Set (First 23 Values)}

See Table~\ref{tab:f:allowed} in \S\ref{sec:frustrated:allowed}.

\subsection{Table 14.3: Example Frustration Values}

See Table~\ref{tab:f:examples} in \S\ref{sec:frustrated:dn}.

\subsection{Table 14.4: Epistemic Status}

\begin{table}[h]
\centering
\caption{Epistemic status of Chapter~\ref{ch:frustrated} claims.}
\label{tab:f:epistemic}
\begin{tabular}{p{5.5cm}cc}
    \toprule
    Claim & Tag & Confidence \\
    \midrule
    Allowed set $S = \{2^a \times 3^b\}$ & [Der] & High \\
    Distance metric $d(n)$ & [Der] & High \\
    Coordination function $n(A) = p A^{1/3}$ & [Cal] & Medium \\
    Prefactor $p \approx 6.1$ & [Cal] & Medium \\
    Linear correction form $\Delta = g \times d(n)$ & [P] & Medium--High \\
    Sign $g < 0$ (frustration $\to$ faster) & [Cal] & High \\
    Magnitude $|g| \approx 1.7$ & [Cal] & Medium \\
    Robustness to confounds & [I] & Medium--High \\
    Forbidden zone [37,47] & [Der] & High \\
    \bottomrule
\end{tabular}
\end{table}

% ============================================================================
\section{Bridge to Chapter~\ref{ch:high_coord}}
\label{sec:frustrated:bridge}
% ============================================================================

Chapter~\ref{ch:high_coord} applies the coordination frustration framework
to make \textbf{predictions} for high-coordination clusters:

\begin{itemize}
    \item Compute $n(A)$ for predicted isotopes
    \item Determine $d(n)$ relative to allowed set
    \item Apply corrected lane with calibrated $g$
    \item Generate testable half-life predictions
\end{itemize}

The forbidden zone [37,47] plays a central role: all high-A clusters
with $A \in [224, 460]$ fall in this forbidden region, experiencing
maximal coordination frustration. \tagP

% ============================================================================
\section{Epistemic Status and Falsifiability}
\label{sec:frustrated:falsifiability}
% ============================================================================

\subsection{Epistemic Status Summary}

See Table~\ref{tab:f:epistemic} for full status.

\subsection{Falsifiability Hooks}

\begin{edcopenbox}[title=Falsifiability Hooks]
    \begin{enumerate}
        \item \textbf{Residual-frustration correlation:}
              If baseline residuals $\Delta$ do \emph{not} correlate with
              $d(n)$ (correlation coefficient $< 0.5$), the EDC correction
              hypothesis fails.

        \item \textbf{Sign of $g$:}
              If future data require $g > 0$ (frustration $\to$ slower release),
              the enhanced-tunneling/preformation interpretation fails.

        \item \textbf{Allowed set violations:}
              If clusters with $n(A) \in S$ (low $d(n)$) show systematic
              deviations comparable to high-$d(n)$ cases, the topological
              origin of $S$ is wrong.

        \item \textbf{Forbidden zone predictions:}
              Clusters falling in the forbidden zone [37,47] must show
              systematically shorter half-lives than baseline. If they match
              baseline, the zone concept is falsified.

        \item \textbf{High-coordination validation:}
              Future measurements of elements 119--120 must fall within
              $\pm 1.5$ dex of corrected-lane predictions. Systematic
              failures invalidate the extrapolation.

        \item \textbf{Alternative $d(n)$ functional forms:}
              If nonlinear corrections (e.g., $d(n)^2$, $\log d(n)$)
              perform significantly better, the linear form is falsified.

        \item \textbf{Prefactor stability:}
              If $p$ must vary systematically with $Z$ or $A$ to maintain
              fit quality, the simple scaling $n(A) = p A^{1/3}$ is inadequate.
    \end{enumerate}
\end{edcopenbox}

% ============================================================================
% Observer-side projection
% ============================================================================

\begin{observerbox}
Observer-side projection note: quoted terms are measurement labels for the
3D/4D projection of the 5D objects defined in this chapter; they do not
imply any conventional mechanism.

\begin{itemize}
    \item \Frustration{} distance $d(n)$ $\leftrightarrow$ structural mismatch proxy (projection label)
    \item Corrected lane residual $\leftrightarrow$ measured time offset (projection label)
\end{itemize}

What the observer records: the frustration correction $g \times d(n)$
improves baseline predictions. Clusters with $n \notin S$ show systematic
deviations that correlate with $d(n)$ at the projection level.
\end{observerbox}

% ============================================================================
\section{Summary and Forward Links}
\label{sec:frustrated:summary}
% ============================================================================

\begin{resultbox}
    \textbf{What this chapter provides:}
    \begin{itemize}
        \item Coordination function $n(A) = p \times A^{1/3}$ [Cal]
        \item Allowed set $S = \{2^a \times 3^b\}$ from Z$_6$ symmetry [Der]
        \item Distance-to-allowed metric $d(n)$ [Der]
        \item Corrected lane: $\Delta \approx g \times d(n)$ with $g < 0$ [P]
        \item Forbidden zone [37,47] and its implications [Der]
        \item 7 falsifiability hooks for testing the framework
    \end{itemize}

    \textbf{What this chapter omits:}
    \begin{itemize}
        \item Numerical value of $g$ (Appendix Q)
        \item Full regression details (Appendix Q)
        \item High-coordination predictions (Chapter~\ref{ch:high_coord})
        \item Conventional terminology mapping (Appendix X)
    \end{itemize}
\end{resultbox}

\bigskip
\noindent
\textbf{Forward references:}
\begin{itemize}
    \item Chapter~\ref{ch:high_coord}: High-coordination predictions using $d(n)$
    \item Appendix~Q: Calibration of $p$ and $g$, robustness analysis
    \item Appendix~X: Mapping to conventional notation
\end{itemize}

\bigskip
\noindent
\textbf{Derivation chain position:}
\begin{equation*}
    \underbrace{\text{Baseline [BL]}}_{\text{Ch.~\ref{ch:barrier_release}}}
    \xrightarrow{+ g \times d(n)}
    \underbrace{\text{Corrected lane}}_{\text{This chapter}}
    \xrightarrow{\text{predictions}}
    \underbrace{\text{High-A}}_{\text{Ch.~\ref{ch:high_coord}}}
\end{equation*}

% ----------------------------------------------------------------------------
% BRIDGE TO NEXT CHAPTER
% ----------------------------------------------------------------------------
\paragraph{Bridge to Chapter~\ref{ch:high_coord}.}
The frustration correction framework is now complete. Chapter~\ref{ch:high_coord}
applies it to the high-coordination regime ($A > 250$)---the superheavy
sector---where baseline errors are largest and frustration effects most
pronounced. The corrected predictions achieve 7$\times$ error reduction
compared to the baseline, providing the strongest quantitative test of
the $\Msix$ topology.

